\subsection{R 中任意符号有限数的可和族}

定理 3. 设 \((x_{t})_{t\in I}\) 是有限实数的一个族；则下述命题等价：

\begin{enumerate}
\def\labelenumi{\alph{enumi})}
\item
  族 \((x_{i})\) 在 \(\mathbf{R}\) 中可和。
\item
  族 \((|x_{i}|)\) 在 \(\mathbf{R}\) 中可和。
\item
  族 \((x_{t})\) 的有限部分和组成的集合在 R 中有界。
\end{enumerate}

设 \(I_1\) 是所有 \(i \in I\) （满足 \(x_i \geq 0\) ）组成的集合，\(I_2\) 是所有 \(i \in I\) （满足 \(x_i < 0\) ）组成的集合。族 \((x_i)_{i \in I}\) {[}相应地，\((|x_i|)_{i \in I}\) {]}可和当且仅当族 \((x_i)_{i \in I_1}\) 与 \((x_i)_{i \in I_2}\) {[}相应地，\((|x_i|)_{i \in I_3}\) 与 \((|x_i|)_{i \in I_4}\) {]}中的每一个都可和（第三章，§ 5，第 3 节，命题 2 与命题 3）。而说 \((x_i)_{i \in I_1}\) 可和、或说 \((|x_i|)_{i \in I_2}\) 可和、或说族 \((x_{i})_{i\in I_1}\) 的有限部分和组成的集合有界（第 1 节，定理 1），这是一回事；把 \(I_1\) 换成 \(I_2\) ，同样的结论也成立。定理随即得出。

定理 3 表明，有限实数族在 R 中的可和性只依赖于其绝对值族的可和性。

我们回忆（第三章，§ 5，第 5 节，命题 6）：若 \((x_{t})\) 与 \((y_{t})\) 是有限实数的两个可和族，则族 \((x_{t} + y_{t})\) 可和，且 \(\sum_{i}(x_{t} + y_{t}) = \sum_{i}x_{t} + \sum_{i}y_{t}\) 。此外，若 \((x_{t})\) 是有限实数的可和族而 \(a\) 是任意有限实数，则族 \((ax_{t})\) 在 \(\mathbf{R}\) 中可和，且我们有 \(\sum_{i}ax_{t} = a\) . \(\sum_{i}x_{t}\) 。

